% Paper for Global WordNet Conference
% Uses ACL style

\documentclass[11pt]{article}
\usepackage[hyperref]{acl2023}
\usepackage{fontspec}
\defaultfontfeatures{Ligatures=TeX}
\setmainfont{TeX Gyre Termes}[Renderer=Harfbuzz,
                              RawFeature={+fallback}]
\setsansfont{TeX Gyre Heros}[Renderer=Harfbuzz,
                             RawFeature={+fallback}]
\setmonofont{TeX Gyre Cursor}[Renderer=Harfbuzz,
                              RawFeature={+fallback} ]
% Script fonts (must be installed on the system)
\newfontfamily\devanagarifont{Noto Serif Devanagari}
\newfontfamily\bengalifont{Noto Serif Bengali}

% Auto-switch by Unicode block
\usepackage{ucharclasses}
\setTransitionsFor{Devanagari}{\devanagarifont}{\normalfont}
\setTransitionsFor{Bengali}{\bengalifont}{\normalfont}

%\usepackage{times}
\usepackage{latexsym}
\usepackage{graphicx}
\usepackage{booktabs}
\usepackage{url}

% For better tables
\usepackage{multirow}

\title{Converting IndoWordNet to the Global WordNet Grid:\\Linking 18 Indian Languages to the Interlingual Index}

\author{Francis Bond \\
  Palacký University Olomouc \\
  \texttt{bond@ieee.org}}

\begin{document}
\maketitle

\begin{abstract}
We present the conversion of IndoWordNet---a linked collection of wordnets for 18 Indian languages---to the standard OMW-LMF (Open Multilingual Wordnet Lexical Markup Framework) format with ILI (Interlingual Index) mappings.
This integration enables these wordnets to interoperate with the Global WordNet Grid, facilitating cross-lingual NLP applications.
We describe the mapping methodology, which handles three types of relationships between IndoWordNet synsets and Princeton WordNet: direct equivalences, hypernymic relationships, and cases where multiple Indian language concepts map to a single English concept.
The conversion preserves IndoWordNet's rich internal semantic relations while adding standardized cross-lingual links.
We release wordnets for Assamese, Bengali, Bodo, Gujarati, Hindi, Kannada, Kashmiri, Konkani, Malayalam, Marathi, Meitei, Nepali, Odia, Punjabi, Sanskrit, Tamil, Telugu, and Urdu.
\end{abstract}

\section{Introduction}

IndoWordNet \cite{bhattacharyya2010indowordnet} is a major lexical resource covering 18 languages of the Indian subcontinent.
Built using the expansion approach with Hindi as the pivot language, it contains over 40,000 synsets per language with rich semantic relations including hypernymy, meronymy, and various language-specific relations.

Despite its size and coverage, IndoWordNet has not been fully integrated into the Global WordNet Grid \cite{bond2016cili}, limiting its interoperability with other multilingual resources.
The primary challenge is that IndoWordNet was built against Princeton WordNet 2.1 \cite{fellbaum1998wordnet}, while the current standard uses the Collaborative Interlingual Index (CILI) based on Princeton WordNet 3.0.

In this paper, we describe our methodology for converting IndoWordNet to the OMW-LMF format with proper ILI mappings.
Our approach addresses three key challenges:
\begin{enumerate}
\item Mapping from PWN 2.1 offsets to PWN 3.0 and then to ILI identifiers
\item Handling cases where the mapping relationship is hypernymic rather than equivalent
\item Preserving lexical distinctions in Indian languages that English conflates
\end{enumerate}

We release converted wordnets for all 18 languages, integrated with the Open Multilingual Wordnet infrastructure.

\section{Background}

\subsection{IndoWordNet}

IndoWordNet \cite{bhattacharyya2010indowordnet} was developed at IIT Bombay using the expansion approach, starting from Hindi WordNet and extending to other Indian languages.
Each language wordnet is linked to the Hindi synsets, which in turn are mapped to Princeton WordNet 2.1.

The resource includes rich semantic relations beyond the standard WordNet taxonomy:
\begin{itemize}
\item Taxonomic: hypernymy, hyponymy, troponymy
\item Meronymy/Holonymy: member, part, portion, substance, location variants
\item Other: also-see, similar, attribute, entailment, causation
\item Language-specific: ability-verb, capability-verb, function-verb, modifies-verb, modifies-noun
\end{itemize}

Access to IndoWordNet is provided through the pyiwn Python API \cite{panjwani2018pyiwn}.

\subsection{The Collaborative Interlingual Index}

The Collaborative Interlingual Index (CILI) \cite{bond2016cili} provides language-neutral concept identifiers that enable wordnets in different languages to be linked.
Each ILI identifier (e.g., \texttt{i105570}) corresponds to a concept, typically grounded in a Princeton WordNet 3.0 synset.

The OMW-LMF format \cite{mccrae2019gwn} is the standard serialization for wordnets in the Global WordNet Association.
It uses XML with ILI attributes on synsets to indicate cross-lingual equivalence.

\subsection{WordNet Version Mapping}

The mapping between Princeton WordNet versions is not straightforward.
Between PWN 2.1 and 3.0, synsets were split, merged, deleted, and renumbered.
We use the official mappings from Princeton \cite{daudert2018pwn21-30} which classify each 2.1-to-3.0 mapping as:
\begin{itemize}
\item \textbf{Equal}: Direct one-to-one correspondence
\item \textbf{Hypernym}: The 3.0 synset is more general
\item \textbf{Hyponym}: The 3.0 synset is more specific
\item Other relationships (auto-hyponym, see-also, etc.)
\end{itemize}

\section{Methodology}

Our conversion pipeline consists of three main stages:
(1) building the ILI mapping from IndoWordNet synsets,
(2) converting synsets and lexical entries to LMF format, and
(3) adding semantic relations with appropriate linking.

\subsection{Building the ILI Mapping}

Each IndoWordNet synset is linked to a PWN 2.1 offset with a relationship type (Direct, Hypernymy, or other).
We chain this through the 2.1-to-3.0 mapping and then look up the ILI from the Open Multilingual Wordnet English resource.

The mapping produces three categories:

\paragraph{Equal mappings}
When the IWN-to-PWN21 relationship is ``Direct'' and the PWN21-to-PWN30 mapping is ``equal,'' we assign the ILI directly to the IndoWordNet synset.
This indicates semantic equivalence.

\paragraph{Hypernym mappings}
When either relationship is hypernymic, the English concept is more general than the Indian language concept.
For example, an IWN synset for a specific type of activity might map to a more general English ``activity'' synset.

\paragraph{Duplicate mappings}
When multiple IWN synsets map to the same ILI, this indicates that the Indian language makes lexical distinctions that English does not.
For example, Hindi distinguishes between अनुभवी (\textit{anubhavī}, ``experienced through skill'') and भुक्तभोगी (\textit{bhuktabhogī}, ``experienced through suffering''), both mapping to the English concept ``experienced.''

Table~\ref{tab:mapping-stats} shows the distribution of mapping types.

\begin{table}[t]
\centering
\begin{tabular}{lr}
\toprule
\textbf{Mapping Type} & \textbf{Count} \\
\midrule
Equal & 12,284 \\
Hypernym & 7,488 \\
Duplicate & 4,482 \\
\midrule
Total mapped & 24,254 \\
\bottomrule
\end{tabular}
\caption{Distribution of ILI mapping types from IndoWordNet to CILI.}
\label{tab:mapping-stats}
\end{table}

\subsection{ILI Assignment Strategy}

Since each ILI should be used only once per wordnet (to maintain the one-to-one correspondence between synsets and concepts), we use the following strategy:

\paragraph{Equal mappings} receive the ILI directly on the synset:
\begin{verbatim}
<Synset id="iwn-hi-s2349-n" 
        ili="i105570">
  <Definition>...</Definition>
</Synset>
\end{verbatim}

\paragraph{Hypernym and duplicate mappings} do not receive a direct ILI.
Instead, we check if another synset already has this ILI through an equal mapping.
If so, we add a hypernym relation to that synset.
If not, we create an ``anchor synset''---a synset with the ILI and English definition but no lemmas---and link the Indian language synset as its hyponym.

\begin{verbatim}
<!-- Anchor synset (English definition) -->
<Synset id="iwn-hi-omw-en-00935500-a" 
        ili="i5127">
  <Definition>having experience...
  </Definition>
</Synset>

<!-- Hindi synset linked as hyponym -->
<Synset id="iwn-hi-s101-a" ili="">
  <Definition>जो भोग चुका हो...
  </Definition>
  <SynsetRelation 
    target="iwn-hi-omw-en-00935500-a" 
    relType="hypernym"/>
</Synset>
\end{verbatim}

For anchor synsets, we also traverse the Princeton WordNet hypernym hierarchy to find the nearest concept that has an equal mapping in IndoWordNet, linking the anchor to this synset.
This maintains taxonomic coherence.

\subsection{Relation Mapping}

IndoWordNet contains rich internal semantic relations that we preserve in the conversion.
Table~\ref{tab:relation-mapping} shows the mapping from IndoWordNet relation types to the standard GWA relation vocabulary.

\begin{table}[t]
\centering
\small
\begin{tabular}{ll}
\toprule
\textbf{IndoWordNet} & \textbf{GWA/OMW-LMF} \\
\midrule
hypernymy & hypernym \\
hyponymy, troponymy & hyponym \\
mero\_member\_collection & mero\_member \\
mero\_component\_object & mero\_part \\
mero\_portion\_mass & mero\_portion \\
mero\_stuff\_object & mero\_substance \\
mero\_place\_area & mero\_location \\
holo\_* & holo\_* (parallel) \\
also\_see & also \\
similar & similar \\
attributes & attribute \\
entailment & entails \\
causative & causes \\
ability\_verb, etc. & other \\
\bottomrule
\end{tabular}
\caption{Mapping of IndoWordNet relations to GWA standard relations.}
\label{tab:relation-mapping}
\end{table}

Language-specific relations without standard equivalents (ability\_verb, capability\_verb, function\_verb, modifies\_verb, modifies\_noun) are mapped to ``other.''

\section{Results}

We successfully converted all 18 IndoWordNet languages to OMW-LMF format.
Table~\ref{tab:language-stats} shows statistics for selected languages.


\input{../build/results.tex}

\begin{table}[t]
\centering
\small
\begin{tabular}{lrrr}
\toprule
\textbf{Language} & \textbf{Synsets} & \textbf{Lemmas} & \textbf{Relations} \\
\midrule
Assamese & 14,958 & 29,432 & 25,815 \\
Bengali & -- & -- & -- \\
Hindi & -- & -- & -- \\
Kannada & -- & -- & -- \\
Malayalam & -- & -- & -- \\
Marathi & -- & -- & -- \\
Tamil & -- & -- & -- \\
Telugu & -- & -- & -- \\
\bottomrule
\end{tabular}
\caption{Statistics for selected IndoWordNet languages after conversion. (To be filled with complete run data.)}
\label{tab:language-stats}
\end{table}

The ILI coverage varies by mapping type:
\begin{itemize}
\item 8,227 synsets (55\%) have direct ILI assignment (equal)
\item 597 synsets (4\%) link to existing equal synsets (hypernym)
\item 3,192 synsets (21\%) link to anchor synsets (duplicate)
\item 2,942 synsets (20\%) have no mapping
\end{itemize}

\subsection{Validation}

We validated the conversion by:
\begin{enumerate}
\item Loading generated LMF files with the standard \texttt{wn} Python library
\item Verifying that known synsets (e.g., ``tree'' पेड़) have correct ILI assignments
\item Checking that hypernym paths are coherent
\item Confirming that internal relations are preserved
\end{enumerate}

\section{Discussion}

\subsection{Lexical Gaps and Language-Specific Concepts}

The duplicate mappings reveal interesting cases where Indian languages make distinctions that English does not lexicalize.
For example, we found 2,068 groups of IndoWordNet synsets that map to the same English concept, including:
\begin{itemize}
\item 1,719 pairs (two Indian concepts per English concept)
\item 288 triples
\item 53 quadruples
\item Higher multiplicities for some concepts
\end{itemize}

These cases represent valuable linguistic data about conceptual distinctions in Indian languages.

\subsection{Missing Mappings}

Approximately 20\% of synsets could not be mapped to ILI.
Analysis shows these fall into several categories:
\begin{itemize}
\item Synsets added to IndoWordNet without PWN links
\item PWN 2.1 synsets that were deleted in PWN 3.0
\item Synsets with non-standard relationship types
\end{itemize}

Future work could address these through manual curation or by proposing new ILI entries.

\subsection{Data Quality Issues}

We encountered some data quality issues in the source IndoWordNet data:
\begin{itemize}
\item Malformed entries where glosses were incorrectly parsed as lemmas
\item Inconsistent relationship annotations
\item A small number of entries with unknown relation types
\end{itemize}

These were handled automatically where possible and flagged for manual review otherwise.

\section{Related Work}

Previous efforts to integrate Indian language wordnets with global resources include the ELRA distribution of Hindi WordNet and partial conversions for specific languages.
Our work is the first comprehensive conversion of all 18 IndoWordNet languages to OMW-LMF with ILI mappings.

The methodology for handling hypernymic and duplicate mappings builds on work in wordnet alignment \cite{daude2003mapping} and the principles established for the Open Multilingual Wordnet \cite{bond2013linking}.

\section{Conclusion}

We have presented the conversion of IndoWordNet to OMW-LMF format with ILI mappings, enabling integration with the Global WordNet Grid.
Our methodology handles the complexities of version mapping and preserves the rich lexical distinctions present in Indian languages.

The converted wordnets are available at [URL to be added] and will be submitted for inclusion in the Open Multilingual Wordnet.

\section*{Acknowledgments}

We thank the IndoWordNet team at IIT Bombay for creating and maintaining this valuable resource, and the developers of the pyiwn library for providing programmatic access.

\bibliography{references}
\bibliographystyle{acl_natbib}

\appendix

\section{Mapping Pipeline}

The conversion pipeline consists of the following scripts:
\begin{enumerate}
\item \texttt{map2ili.py}: Builds ILI mappings from IndoWordNet via PWN 2.1 $\rightarrow$ 3.0 $\rightarrow$ ILI chain
\item \texttt{iwn2omw.py}: Converts to OMW-LMF format with relations
\item \texttt{test\_lmf.py}: Validates generated wordnets
\end{enumerate}

All code is available at [GitHub URL].

\end{document}